\section{Formal Setup}\label{sec:formal-setup}

This section restates the definitions needed for the paper to stand on its own. The emphasis is on sufficiency, relevance, and the quotient structure of the optimizer map.

\subsection{Decision Problems and Sufficiency}

\begin{definition}[Decision Problem]\label{def:decision-problem}
A decision problem with coordinate structure is a tuple $\mathcal{D}=(A,X_1,\ldots,X_n,U)$ where $A$ is a finite action set, $S=X_1 \times \cdots \times X_n$ is the state space, and $U : A \times S \to \mathbb{Q}$ is the utility function.
\end{definition}

\begin{definition}[Projection]\label{def:projection}
For state $s=(s_1,\ldots,s_n)\in S$ and coordinate set $I \subseteq \{1,\ldots,n\}$, write $s_I$ for the projection of $s$ onto the coordinates in $I$.
\end{definition}

\begin{definition}[Optimizer Map]\label{def:optimizer}
The optimizer map is
\[
\Opt(s) := \arg\max_{a \in A} U(a,s).
\]
\end{definition}

\begin{definition}[Sufficient Coordinate Set]\label{def:sufficient}
Coordinate set $I$ is sufficient if
\[
\forall s,s' \in S:\quad s_I=s'_I \implies \Opt(s)=\Opt(s').
\]
\end{definition}

\begin{definition}[Minimal Sufficient Set]\label{def:minimal-sufficient}
A sufficient set is minimal if no proper subset is sufficient.
\end{definition}

\begin{definition}[Relevant Coordinate]\label{def:relevant}
Coordinate $i$ is relevant if there exist states that differ only at coordinate $i$ and induce different optimal-action sets:
\[
\exists s,s' \in S:\;
\Big(\forall j \neq i,\; s_j=s'_j\Big)
\ \wedge\
\Opt(s)\neq\Opt(s').
\]
\end{definition}

\begin{proposition}[Minimal Sufficient Sets Are Relevant Coordinates]\label{prop:minimal-relevant-equiv}
Every minimal sufficient set is exactly the relevant-coordinate set.
\end{proposition}

\begin{proof}
\textbf{Minimal sufficient sets contain only relevant coordinates.}
Let $I$ be minimal sufficient and let $i\in I$. Suppose toward contradiction that $i$ is not relevant. Then whenever two states agree on all coordinates except possibly $i$, they still have the same optimizer value. In particular, once the coordinates in $I\setminus\{i\}$ are fixed, changing coordinate $i$ cannot change $\Opt$.

Now take any two states $s,s'$ with $s_{I\setminus\{i\}}=s'_{I\setminus\{i\}}$. Keeping the other coordinates fixed, we can vary coordinate $i$ without affecting $\Opt$, so the value of $\Opt$ is already determined by the coordinates in $I\setminus\{i\}$. Thus $I\setminus\{i\}$ is still sufficient, contradicting minimality of $I$. Therefore every coordinate in a minimal sufficient set is relevant.

\textbf{Every relevant coordinate belongs to every sufficient set.}
Let $i$ be relevant. By definition, there exist states $s,s'$ differing only at coordinate $i$ such that $\Opt(s)\neq\Opt(s')$. If a coordinate set $I$ omitted $i$, then these two states would agree on all coordinates of $I$, since they differ nowhere else. Sufficiency of $I$ would then force $\Opt(s)=\Opt(s')$, contradiction. So every sufficient set must contain $i$.

Combining the two parts, every minimal sufficient set is contained in the relevant-coordinate set, and the relevant-coordinate set is contained in every sufficient set, hence in every minimal sufficient set. Therefore every minimal sufficient set is exactly the relevant-coordinate set.
\end{proof}

\begin{definition}[Structural Rank]\label{def:srank}
The structural rank of $\mathcal{D}$ is the cardinality of the relevant-coordinate support:
\[
\mathrm{srank}(\mathcal{D}) :=
\left|\{i \in \{1,\ldots,n\} : i \text{ is relevant}\}\right|.
\]
\end{definition}

\begin{proposition}[Structural Rank as Relevant-Support Size]\label{prop:srank-support}
\[
\mathrm{srank}(\mathcal{D})
=
\left|\{i \in \{1,\ldots,n\} : i \text{ is relevant}\}\right|.
\]
Hence $\mathrm{srank}(\mathcal{D}) \le n$, and $\mathrm{srank}(\mathcal{D})=0$ exactly when no coordinate is relevant.
\leanmeta{\LHrng{SK}{1}{3}.}
\end{proposition}

\begin{proof}
The equality is the definition of structural rank. The upper bound is immediate because the relevance support is a subset of $\{1,\dots,n\}$. The zero case is equivalent to the relevance support being empty.
\end{proof}

This quantity is the central invariant of the paper. The remaining sections ask which different mathematical summaries of a decision problem are in fact measuring this same support-size core.

\subsection{Optimizer Quotient and Coimage Structure}

\begin{definition}[Decision Equivalence]\label{def:decision-equiv}
States $s,s' \in S$ are decision-equivalent if $\Opt(s)=\Opt(s')$.
\end{definition}

\begin{definition}[Decision-Quotient Score]\label{def:decision-quotient}
For coordinate set $I$ and state $s$, define
\[
\DQ_I(s) =
\frac{|\{a \in A : a \in \Opt(s') \text{ for some } s' \text{ with } s'_I=s_I\}|}{|A|}.
\]
\end{definition}

\begin{proposition}[Sufficiency Characterization]\label{prop:sufficiency-char}
Coordinate set $I$ is sufficient if and only if $\DQ_I(s)=|\Opt(s)|/|A|$ for every state $s$.
\end{proposition}

\begin{proof}
\textbf{If $I$ is sufficient.} Fix $s$. For any $s'$ with $s'_I=s_I$, sufficiency gives $\Opt(s')=\Opt(s)$. Hence
\[
\{a : a\in \Opt(s') \text{ for some } s' \text{ with } s'_I=s_I\}=\Opt(s),
\]
so $\DQ_I(s)=|\Opt(s)|/|A|$.

\textbf{If the displayed equality holds for all $s$.} Suppose $I$ were not sufficient. Then there exist $s,s'$ with $s_I=s'_I$ but $\Opt(s)\neq\Opt(s')$. Pick $a\in \Opt(s')\setminus \Opt(s)$ (or symmetrically). This action belongs to the union in the numerator of $\DQ_I(s)$ but not to $\Opt(s)$, so the numerator is strictly larger than $|\Opt(s)|$, contradiction.
\end{proof}

\begin{proposition}[Optimizer Quotient as Coimage/Image Factorization]\label{prop:optimizer-coimage}
The quotient of $S$ by decision equivalence is canonically equivalent to $\operatorname{range}(\Opt)$. Equivalently, the optimizer quotient is the coimage of the optimizer map, identified with its image in \textbf{Set}.
\end{proposition}

\begin{proof}
Define
\[
\Phi: S/{\sim}\ \to\ \operatorname{range}(\Opt),\qquad \Phi([s])=\Opt(s).
\]
This is well-defined because $s\sim s'$ implies $\Opt(s)=\Opt(s')$. It is surjective by definition of range. It is injective because $\Phi([s])=\Phi([s'])$ means $\Opt(s)=\Opt(s')$, hence $[s]=[s']$. Therefore $\Phi$ is a bijection.
\end{proof}

\begin{theorem}[Universal Property of the Optimizer Quotient]\label{thm:quotient-universal}
Let $Q=S/{\sim}$ be the optimizer quotient, where $s \sim s'$ if and only if $\Opt(s)=\Opt(s')$, and let $\pi:S\to Q$ be the quotient map. For any surjective abstraction $\phi:S\to T$ that preserves the optimizer map in the sense that
\[
\phi(s)=\phi(s') \implies \Opt(s)=\Opt(s'),
\]
there exists a unique map $\psi:T\to Q$ such that $\pi=\psi\circ\phi$.
\end{theorem}

\begin{proof}
Because $\phi$ preserves optimizer values on fibers, we have
\[
\phi(s)=\phi(s') \implies s\sim s'.
\]
So each $\phi$-fiber is contained in a $\sim$-class.

For each $t\in T$, choose any $s$ with $\phi(s)=t$ (possible since $\phi$ is surjective), and define
\[
\psi(t):=\pi(s)\in Q.
\]
This is well-defined: if $s_1,s_2$ are two preimages of $t$, then $\phi(s_1)=\phi(s_2)$, hence $s_1\sim s_2$, hence $\pi(s_1)=\pi(s_2)$.

Now for any $s\in S$,
\[
(\psi\circ\phi)(s)=\psi(\phi(s))=\pi(s),
\]
so $\pi=\psi\circ\phi$.

For uniqueness, let $\psi'$ also satisfy $\pi=\psi'\circ\phi$. For any $t\in T$, choose $s$ with $\phi(s)=t$. Then
\[
\psi'(t)=\psi'(\phi(s))=\pi(s)=\psi(\phi(s))=\psi(t),
\]
so $\psi'=\psi$.
\end{proof}

\begin{proposition}[Optimizer Quotient Entropy Depends Only on the Image]\label{prop:optimizer-entropy-image}
If $\mathrm{numOptClasses}(\mathcal{D})$ denotes the number of distinct values attained by $\Opt$, then
\[
H(\mathcal{D}) = \log_2(\mathrm{numOptClasses}(\mathcal{D}))
=
\log_2(|\operatorname{range}(\Opt)|).
\]
\end{proposition}

\begin{proof}
By Proposition~\ref{prop:optimizer-coimage}, quotient classes are in bijection with $\operatorname{range}(\Opt)$, so their cardinalities agree:
\[
\mathrm{numOptClasses}(\mathcal{D})=|\operatorname{range}(\Opt)|.
\]
By definition, $H(\mathcal{D})$ is the base-$2$ entropy of the counting-normalized quotient distribution; on a finite uniform support of size $m$, entropy is $\log_2 m$. Substituting $m=\mathrm{numOptClasses}(\mathcal{D})$ yields the claim.
\end{proof}

\begin{remark}[Decision Questions vs Decision Problems]\label{rem:question-vs-problem}
The predicate ``is $I$ sufficient for $\mathcal{D}$?'' is encoding-independent. Complexity classes apply only after choosing how $\mathcal{D}$ and $I$ are represented. This paper uses the same shared definitions, but studies their quotient and convergence consequences rather than their class-theoretic complexity.
\end{remark}

The point of Theorem~\ref{thm:quotient-universal} is organizational rather than terminological. In \textbf{Set}, the optimizer quotient is a familiar coimage/image construction. Its role here is to provide a common comparison object for entropy, support counting, Bayesian updating, and lossless summaries.
